\cotefile{34}
\scanpage{3}
\mspage{1}
\uline{Exposé I}\qquad\deleted{\uline{Action du groupe de monodromie :}}
\uline{Généralités, \inserted{le} th. de monodromie par voie arithmétique.}

0. \uline{Introduction}
\figure{\begin{tikzpicture}[x=1cm,y=1cm]\draw (0,1)..controls(1.2,1.1) and(1.1,0)..(2.2,0)--(12,0.15);\end{tikzpicture}}

0.1. \uline{\unsure{Notations}}. $S=\operatorname{Spec}V$ un trait, $s$ pt fermé, $\eta$ point générique
\[
\mathfrak m=\pi V\inserted{\text{ ``}V\text{''}},\quad k=k(s)\inserted{\text{ ``}V/\mathfrak m\text{''}},\quad K=k(\eta);\qquad p=\text{exp. car. de }k(s).
\]
Le plus souvent, on supposera $V$ hensélien, et \unsure{$S$|même} strictement \unsure{locale|local} i.e. $k$ sép. clos;

\deleted{\unsure{$K$}} $\bar K$ une clôture séparable de $K$, $\bar\eta=\operatorname{Spec}\bar K$;
\[
\pi=\operatorname{Gal}(\bar K/K).
\]
$\ell$ nombre premier. \deleted{Sauf mention du contraire} \deleted{Dans la plus grande partie \unsure{voir}} \inserted{\unsure{Pour toute la première partie du séminaire}, on supposera}
\deleted{L'hyp. de}
\[
\ell\neq\operatorname{car}k
\]

0.2. L'objet du séminaire est l'étude de certaines représentations $\ell$-adiques de $\pi$, provenant de diverses situations géométriques. Une repr. $\ell$-adique de $\pi$ est un \deleted{\unsure{module}} hom. continu
\[
(*)\qquad\pi\longrightarrow\operatorname{Aut}_{\mathbb Z_\ell}(E),
\]
où $E$ est un module de t.f. sur $\mathbb Z_\ell$. On pourra aussi regarder
\[
\pi\longrightarrow\operatorname{Aut}_{\mathbb Q_\ell}(E)
\]
$E$ module de d.f. sur $\mathbb Q_\ell$.

\uline{Exemple} 0.3. $X_K$ schéma algébrique sur $K$, \inserted{$i$ entier.} On suppose que $H^i(X_{\bar K},\mathbb Z/\ell\mathbb Z)$ est un groupe fini, de sorte que
\[
E=H^i(X_{\bar K},\mathbb Z_\ell)\overset{\mathrm{dfn}}=\underset{\nu}{\unsure{\uline{\lim}|\varprojlim}}H^i(X_{\bar K},\mathbb Z/\ell^\nu\mathbb Z)
\]
est un module de t.f. sur $\mathbb Z_\ell$; \deleted{Cette hyp. est} \deleted{\unsure{est}} vérifiée dans les trois cas suivants :

a) On dispose de résolution des sing. \inserted{sous forme de Hironaka (\quad)} des schémas alg. de dim $\leq\dim X$ sur $\bar K$, -- p.ex. $K$ de car. nulle

b) $X$ propre sur $K$.

c) $X$ lisse sur $K$.
\margin{2}
